\documentclass[11pt,a4paper]{article}

% ---------- PACOTES ----------
\usepackage[utf8]{inputenc}
\usepackage[T1]{fontenc}
\usepackage[brazilian]{babel}
\usepackage{lmodern}
\usepackage[a4paper,margin=2.3cm,top=2.6cm,bottom=2.8cm,headheight=22pt]{geometry}
\usepackage{xcolor}
\usepackage{colortbl}
\usepackage{booktabs}
\usepackage{longtable}
\usepackage{array}
\usepackage{tikz}
\usepackage{titlesec}
\usepackage{enumitem}
\usepackage{fancyhdr}
\usepackage{microtype}
\usepackage{hyperref}
\usepackage{textcomp}

% ---------- PALETA ----------
\definecolor{voltnavy}{HTML}{10151B}
\definecolor{voltcyan}{HTML}{1F6F8F}
\definecolor{voltcyanlight}{HTML}{5BC2E0}
\definecolor{voltcopper}{HTML}{B5652F}
\definecolor{voltamber}{HTML}{9A640C}
\definecolor{voltamberbg}{HTML}{F3E6CD}
\definecolor{voltok}{HTML}{2E6B4C}
\definecolor{voltokbg}{HTML}{DFEEE6}
\definecolor{voltpending}{HTML}{5B5F82}
\definecolor{voltgray}{HTML}{4D5A63}
\definecolor{voltrule}{HTML}{CDD4D9}
\definecolor{voltbandbg}{HTML}{EEF1F2}

\hypersetup{
  colorlinks=true,
  linkcolor=voltcyan,
  urlcolor=voltcyan,
  citecolor=voltcyan,
  pdftitle={SkyVolt - Documento de Requisitos, Arquitetura e BOM},
  pdfsubject={Drone medidor e fiscalizador de ambientes de alta e baixa energia}
}

% ---------- TIPOGRAFIA DE SEÇÕES ----------
\titleformat{\section}
  {\Large\bfseries\color{voltcyan}}
  {\thesection}{0.6em}{}
\titlespacing*{\section}{0pt}{1.6em}{0.9em}

\titleformat{\subsection}
  {\large\bfseries\color{voltcopper}}
  {\thesubsection}{0.6em}{}
\titlespacing*{\subsection}{0pt}{1.2em}{0.6em}

\renewcommand{\contentsname}{Sumário}

% ---------- CABEÇALHO/RODAPÉ ----------
\pagestyle{fancy}
\fancyhf{}
\renewcommand{\headrulewidth}{0.4pt}
\renewcommand{\headrule}{\hbox to\headwidth{\color{voltrule}\leaders\hrule height \headrulewidth\hfill}}
\fancyhead[L]{\footnotesize\ttfamily\color{voltgray} SkyVolt}
\fancyhead[R]{\footnotesize\ttfamily\color{voltgray} Requisitos \& Arquitetura}
\fancyfoot[C]{\footnotesize\ttfamily\color{voltgray} \thepage}

% ---------- COMANDOS AUXILIARES ----------
\newcommand{\specline}[2]{{\ttfamily\footnotesize\color{voltcyanlight}#1}\hfill{\normalsize\color{white}#2}\\[0.4em]}
\newcommand{\statusok}{\textcolor{voltok}{\textbf{decidido}}}
\newcommand{\statuspend}{\textcolor{voltpending}{\textbf{pendente}}}
\newcommand{\statushaz}{\textcolor{voltamber}{\textbf{limitação}}}
\newcommand{\pesq}{\textsuperscript{\textdagger}}

% caixa de destaque sem depender de tcolorbox
\newcommand{\aviso}[2]{%
  \par\noindent\colorbox{voltamberbg}{%
    \begin{minipage}{\dimexpr\linewidth-2\fboxsep\relax}
      \vspace{0.3em}
      \textcolor{voltamber}{\textbf{#1}} #2
      \vspace{0.3em}
    \end{minipage}%
  }\par\vspace{1em}
}

\newcommand{\destaque}[2]{%
  \par\noindent\colorbox{voltbandbg}{%
    \begin{minipage}{\dimexpr\linewidth-2\fboxsep\relax}
      \vspace{0.3em}
      \textcolor{voltcyan}{\textbf{#1}} #2
      \vspace{0.3em}
    \end{minipage}%
  }\par\vspace{1em}
}

\setlength{\parskip}{0.55em}
\setlength{\parindent}{0pt}

\begin{document}

% =====================================================
% CAPA
% =====================================================
\begin{titlepage}
\pagecolor{voltnavy}
\color{white}
\thispagestyle{empty}

\begin{center}
\vspace*{0.6cm}
{\ttfamily\small\color{voltcyanlight} DOCUMENTO DE REQUISITOS, ARQUITETURA \& BOM \;\textbullet\; REV.\ 1.3}

\vspace{1.6cm}

\begin{tikzpicture}[scale=1.0]
  \foreach \a in {0,60,...,300} {
    \draw[voltcyanlight!55!white, line width=0.7pt] (0,0) -- (\a:2.15);
    \filldraw[voltnavy, draw=voltcyanlight!75!white, line width=0.8pt] (\a:2.15) circle (3.1pt);
  }
  \filldraw[voltnavy, draw=white, line width=1pt] (0,0) circle (7pt);

  % efetuador — plugue (lateral, ciano)
  \draw[voltcyanlight, line width=1.1pt] (0,0) -- (0:1.35);
  \filldraw[voltcyanlight] (0.1,1.2) rectangle (0.42,1.5);
  \node[anchor=west, voltcyanlight, font=\ttfamily\tiny] at (0.9,2.3) {plugue 3 pinos};
  \draw[voltcyanlight, line width=0.4pt] (0.42,1.4) -- (0.85,2.25);

  % efetuador — garra de campo (âmbar)
  \draw[voltamber!70!voltcyanlight, line width=1.1pt] (0,0) -- (215:1.3);
  \draw[voltamber!60!white, line width=1.3pt] (215:1.3) circle (4.4pt);
  \node[anchor=east, voltamber!60!white, font=\ttfamily\tiny] at (-2.0,-1.55) {garra de campo};
  \draw[voltamber!60!white, line width=0.4pt] (215:1.3) -- (-2.0,-1.5);

  % efetuador — garra de contato (cobre)
  \draw[voltcopper!85!white, line width=1.1pt] (0,0) -- (145:1.3);
  \filldraw[voltcopper!85!white] (145:1.3) circle (4.4pt);
  \node[anchor=east, voltcopper!85!white, font=\ttfamily\tiny] at (-2.0,1.55) {garra de contato};
  \draw[voltcopper!85!white, line width=0.4pt] (145:1.3) -- (-2.0,1.5);
\end{tikzpicture}

\vspace{0.9cm}

{\ttfamily\footnotesize\color{voltcyanlight} NOME DE PROJETO}\\[0.35em]
{\fontsize{42}{46}\selectfont\bfseries SkyVolt}\\[0.55em]
{\normalsize\itshape\color{white!70!voltnavy} Drone medidor e fiscalizador de ambientes de alta e baixa energia}

\vspace{1.3cm}

\begin{minipage}{10.6cm}
\specline{CURSO}{Desenvolvimento de Aplicações Desktop}
\specline{DATA}{14/08/2026}
\specline{REVISÃO}{1.3}
\specline{IMPRESSÃO 3D}{Fornecida pela faculdade — sem custo ao projeto}
\specline{TETO DE ORÇAMENTO (EQUIPAMENTO)}{R\$~2.000,00}
\specline{GASTO MÉDIO ESTIMADO (BOM)}{R\$~1.903,80}
\specline{ELABORADO POR}{\_\_\_\_\_\_\_\_\_\_\_\_\_\_\_\_\_\_\_\_\_\_\_\_\_\_\_\_}
\end{minipage}

\vfill
{\ttfamily\tiny\color{voltcyanlight!70!white} vista superior esquemática — três efetuadores em posições distintas do drone}
\vspace{0.4cm}
\end{center}
\end{titlepage}

\pagecolor{white}
\color{black}

% =====================================================
% SUMÁRIO
% =====================================================
\thispagestyle{empty}
\tableofcontents
\newpage

% =====================================================
% 1. VISÃO GERAL
% =====================================================
\section{Visão geral do projeto}

O SkyVolt é um hexacóptero de estrutura impressa em 3D (impressora fornecida pela faculdade, sem custo ao projeto) cuja função principal é medir tensão e corrente elétrica em pontos de um ambiente controlado, classificando automaticamente o tipo de circuito e reportando os dados a um aplicativo desktop de controle e monitoramento.

\destaque{Princípio orientador de escopo:}{a medição é a função \emph{core} do projeto; a plataforma voadora é secundária. Sempre que o orçamento ou o prazo exigirem corte de escopo, a prioridade recai sobre manter a função de medição íntegra, simplificando o drone antes de simplificar o sensoriamento.}

\textbf{Papel do aplicativo desktop} (entregável central da disciplina): estação de controle e monitoramento, recebendo telemetria em tempo real via LoRa, exibindo tensão, corrente, temperatura e classificação num painel, mantendo histórico local de medições e emitindo alertas configuráveis.

\textbf{Faixa elétrica do escopo:} circuitos AC de até 220V (padrão residencial/tomada) e circuitos DC sempre $\leq$100V (restrição de projeto). Ambientes acima de 220V estão fora de escopo pelo risco de arco elétrico.

% =====================================================
% 2. REQUISITOS FUNCIONAIS
% =====================================================
\section{Requisitos funcionais}

\begin{longtable}{@{}p{1.6cm}p{12.6cm}@{}}
\toprule
\textbf{ID} & \textbf{Requisito} \\
\midrule
\endhead
\bottomrule
\endfoot
RF01 & Medir tensão AC/DC entre 5V e 220V, conforme o modo de contato disponível. \\
RF02 & Medir corrente elétrica, por contato direto ou por indução (clamp sem contato). \\
RF03 & Estimar potência dissipada (Lei de Joule) a partir de leitura térmica. \\
RF04 & Transmitir telemetria de sensores (tensão, corrente, classificação AC/DC, modo utilizado) via LoRa até a estação base. \\
RF05 & Exibir os dados recebidos em tempo real no aplicativo desktop (dashboard). \\
RF06 & Registrar histórico de medições localmente (SD/flash) e sincronizar com o app quando o link permitir. \\
RF07 & Posicionar um entre três efetuadores — garra de campo, garra de contato ou plugue de 3 pinos — no ponto de medição, conforme o modo classificado. \\
RF08 & Emitir alerta no aplicativo quando valores excederem limites configuráveis. \\
RF09 & Transmitir vídeo de navegação (FPV) e comandos de voo via WiFi, com baixa latência. \\
RF10 & Classificar automaticamente o alvo (conector padrão / campo AC $\geq$100V / contato $<$100V ou DC) por reconhecimento sem contato, a uma distância de stand-off calibrada, sem informação prévia do operador. \\
RF11 & Executar rotina autônoma de segurança ao detectar qualquer falha, navegando até uma zona segura pré-estabelecida por marcador visual. \\
\end{longtable}

% =====================================================
% 3. RNF
% =====================================================
\section{Requisitos não-funcionais e restrições}

\begin{longtable}{@{}p{1.9cm}p{12.3cm}@{}}
\toprule
\textbf{ID} & \textbf{Restrição} \\
\midrule
\endhead
\bottomrule
\endfoot
RNF01 & Teto de orçamento: R\$~2.000 para equipamento. Impressão 3D fornecida pela faculdade (acordo com o coordenador), sem custo ao projeto. \\
RNF02 & Isolamento galvânico entre o circuito medido (até 220V) e a eletrônica de controle; hardware do modo contato dimensionado para tolerar até 220V mesmo em caso de erro de classificação (defesa em profundidade). \\
RNF03 & Domínios de energia separados — propulsão, eletrônica/medição e supervisor de segurança — para evitar falha compartilhada. \\
RNF04 & Tolerância de posicionamento de $\pm$1,5cm por eixo na abordagem manual, compensada por folga mecânica e funil de captura. \\
RNF05 & Autonomia mínima de voo de 8 a 12 minutos por missão. \\
RNF06 & Escopo elétrico fixo: DC sempre $\leq$100V; AC até 220V. \\
RNF07 & Redundância mínima de propulsão: tolerância a falha de 1 motor entre 6 (hexacóptero). \\
\end{longtable}

% =====================================================
% 4. ARQUITETURA
% =====================================================
\section{Arquitetura do sistema}

\subsection{Estrutura e propulsão \normalfont\statusok}

Estrutura impressa em impressora 3D fornecida pela faculdade (acordo com o coordenador, sem custo ao orçamento). Configuração \textbf{hexacóptero} (6 motores), escolhida como meio-termo entre custo e redundância mecânica: tolera a perda de 1 motor sem queda imediata, ao contrário de um quadricóptero, e custa/pesa menos que o octocóptero originalmente cogitado.

\textbf{Motores brushless RS1606 3300KV} (hélice 3", ESC brushless dedicado por motor) — decisão revisada na Rev.~1.2 a partir do fechamento da verificação de empuxo abaixo. A escolha original por motor de escova (Rev.~1.1) foi descartada: motor de escova compatível com o porte do frame não sustenta o peso estimado da aeronave. Motor brushless exige comutação eletrônica trifásica — um driver ponte-H simples (ex.\ L293D), adequado para motor de escova, \emph{não} serve para acionar motor brushless; por isso cada motor tem um ESC dedicado no BOM (Seção~6.1).

\destaque{Verificação de empuxo (thrust-to-weight) — fechada:}{peso bruto máximo estimado (cenário pessimista, todos os componentes) $\approx$1.110g — atualizado na Rev.~1.3 com o peso real da estrutura, obtido do CAD paramétrico (FreeCAD) em vez de estimativa de densidade: 264g para hub + 6 braços + bandeja de eletrônica (100\% infill PETG). Motor RS1606~3300KV entrega até 875g de empuxo por motor com hélice de 4" (dado de fabricante, 23,7A/398W a 4S) — TWR$\approx$4,7. Com hélice de 3" (empuxo reduzido, estimado por teoria de disco atuador na ausência de bancada publicada para esse par), TWR estimado $\approx$2,7--3,5, ainda acima do mínimo seguro de 1,5--2,0. A hélice de 3" foi escolhida no lugar da de 4" para também trazer a autonomia para dentro da janela de RNF05 sem alterar a bateria — ver Seção~4.4.}

\textbf{Estrutura modular (Rev.~1.3):} o chassi foi redesenhado em 3 peças impressas separadamente — \textbf{hub} (129,6$\times$121,2mm), \textbf{braço} (154$\times$26mm, imprimir $\times$6) e \textbf{bandeja de eletrônica} (150$\times$80mm) — todas dentro da mesa de impressão de 220$\times$220mm do equipamento disponível (o frame original de peça única, 330$\times$290mm, não cabia). Braço e hub se unem por encaixe lingueta/rasgo (tenon/mortise) reforçado com 2 parafusos M3 por junta (12 no total), mais resistente à flexão do que um encaixe topo-a-topo só parafusado. Pad do motor dimensionado para o RS1606~3300KV: padrão de furação quadrado 12$\times$12mm, 4$\times$M2 — \emph{padrão típico da classe 1606, não confirmado contra o datasheet oficial do fabricante (ver Seção~8)}.

\subsection{Sistema de medição — função principal \normalfont\statusok}

Três efetuadores especializados, cada um resolvendo um modo de medição diferente:

\begin{itemize}[itemsep=2pt]
  \item \textbf{Garra de campo} (servo, retrátil) — medição sem contato, para AC $\geq$100V.
  \item \textbf{Garra de contato} (servo) — contato direto, para $<$100V ou DC.
  \item \textbf{Plugue de 3 pinos} (padrão NBR~14136), montado na lateral do drone, com funil/guia de captura impresso e ponta modular para adaptadores de outros tipos de tomada.
\end{itemize}

A escolha do efetuador segue uma triagem em três etapas: aproximação até uma distância de stand-off calibrada, reconhecimento sem contato (sensor de campo EMF + câmera térmica) e classificação — tudo decidido pelo próprio sistema, sem informação prévia do operador sobre o circuito, conforme resumido na tabela a seguir.

\textbf{Fixação física (Rev.~1.3):} os 6 braços do frame já saem impressos com base de fixação para servo SG90 (furação M2, espaçamento 32,5mm, padrão de mercado) — os 3 braços efetuadores usam essa base para a garra/plugue; os outros 3 ficam com a base livre, sem custo extra relevante de peso ou impressão. Carcaça específica de cada efetuador (garra de campo, garra de contato, funil do plugue NBR~14136) ainda não foi desenhada — ver Seção~8.

\textbf{Árvore de triagem do sistema de medição}\\[0.4em]
\begin{longtable}{@{}p{4.6cm}p{4.6cm}p{5cm}@{}}
\toprule
\textbf{Condição} & \textbf{Modo escolhido} & \textbf{Precisão obtida} \\
\midrule
\endfirsthead
\toprule
\textbf{Condição} & \textbf{Modo escolhido} & \textbf{Precisão obtida} \\
\midrule
\endhead
\bottomrule
\endfoot
Conector padrão visível & Modo A — plugue 3 pinos & Tensão e corrente exatas, seguro em qualquer tensão do escopo \\
Sinal AC no campo, $\geq$100V & Modo B — campo (clamp) & Corrente exata; tensão em faixa (limitação física) \\
Sinal AC no campo, $<$100V & Modo B — contato & Tensão e corrente exatas, isolado até 220V \\
Sem sinal AC (DC, $\leq$100V garantido) & Modo B — contato & Tensão e corrente exatas \\
\end{longtable}

\subsection{Comunicação e redundância \normalfont\statusok}

Dois rádios com papéis distintos, mapeados a dois microcontroladores independentes:

\begin{itemize}[itemsep=2pt]
  \item \textbf{WiFi (ESP32~\#1):} vídeo de navegação (FPV) + comandos de voo — baixa latência.
  \item \textbf{LoRa (ESP32~\#2):} telemetria de sensores — pacotes pequenos, atua também como \textbf{supervisor de segurança}.
\end{itemize}

Os dois ESP32 trocam um heartbeat por fio (independente de rádio). O ESP32~\#2 tem câmera própria dedicada e fonte de energia independente — nenhuma dependência compartilhada com o ESP32~\#1 que pudesse derrubar os dois juntos.

\textbf{Máquina de failsafe:} qualquer problema detectado (perda de link, bateria crítica, falha de sensor/IMU, efetuador travado) aciona a mesma rotina — se engajado com o alvo, desengate controlado e lento primeiro; em seguida, navegação autônoma via marcador visual até a zona segura pré-estabelecida, pouso com motores desarmados e efetuador inerte.

\subsection{Energia \normalfont\statusok}

\begin{longtable}{@{}p{3cm}p{4.3cm}p{3cm}p{5cm}@{}}
\toprule
\textbf{Domínio} & \textbf{Química/fonte} & \textbf{Proteção} & \textbf{Observação} \\
\midrule
\endhead
\bottomrule
\endfoot
Propulsão & Li-ion alta descarga (18650/21700) & BMS obrigatório & Mais seguro que LiPo em queda perto de circuito energizado \\
Eletrônica/medição & Domínio próprio, separado da propulsão & Isolamento do sag de tensão & Evita brownout do ESP32 durante o contato \\
Supervisor (ESP32~\#2) & Fonte independente dedicada & — & Sobrevive a falha da bateria de eletrônica \\
\end{longtable}

Autonomia alvo por missão: 8 a 12 minutos (decolagem $\rightarrow$ aproximação $\rightarrow$ medição $\rightarrow$ retorno $\rightarrow$ pouso).

\destaque{Estimativa de autonomia (atualizada na Rev.~1.3 com o peso real da estrutura):}{motor RS1606~3300KV + hélice~3", bateria de propulsão mantida sem alteração (4$\times$ células 21700 alta descarga, 4,2Ah, 4S1P, $\approx$48,4Wh utilizáveis). Potência em hover estimada em $\approx$277W (era 248W na Rev.~1.2, antes do peso real do CAD) $\rightarrow$ autonomia $\approx$10,5 minutos, ainda dentro da janela de RNF05 mas com menos folga que antes. Estimativa apoiada em extrapolação de curva de eficiência de motor similar (sem bancada real do par motor+hélice escolhido) — ver pendência na Seção~8.}

\subsection{Iluminação \normalfont\statusok}

LED difuso, para a câmera normal de navegação (FPV) em ambientes escuros — difuso em vez de focado para evitar reflexo em contatos metálicos. Controle manual, ligado/desligado pelo piloto.

\subsection{PCBs \normalfont\statusok}

Placas perfuradas de 1 camada, soldadas à mão, divididas por domínio funcional: potência/propulsão, ESP32~\#1, ESP32~\#2, sensores de medição (duas metades fisicamente separadas — lado perigoso e lado seguro, isolados por optoacoplador certificado) e adaptador do Modo~A. Sem plano de terra, a separação física da placa de potência e o aterramento em estrela ganham ainda mais importância para reduzir ruído de chaveamento contaminando a leitura.

\textbf{Fixação física (Rev.~1.3):} bandeja impressa dedicada (150$\times$80mm), aparafusada nos 4 furos de coluna do hub por espaçadores/standoffs — carrega os 2 ESP32-S3 em trilhos dimensionados para o footprint oficial do DevKitC-1 (70$\times$28mm) e os 2 módulos LoRa em baia genérica com abraçadeira (footprint do módulo real ainda não confirmado — ver Seção~8).

% =====================================================
% 5. RISCOS
% =====================================================
\section{Riscos e limitações conhecidas}

\begin{longtable}{@{}p{0.7cm}p{9.6cm}p{3.4cm}@{}}
\toprule
\textbf{\#} & \textbf{Risco / limitação} & \textbf{Status} \\
\midrule
\endhead
\bottomrule
\endfoot
1 & Cenário máximo (R\$~2.674,70) ainda ultrapassa o teto de R\$~2.000 caso cotações finais batam no topo da faixa. & \statushaz \\
2 & Redundância de motor reduzida (hexacóptero vs.\ octocóptero original). & troca aceita \\
3 & Empuxo e autonomia do par motor RS1606~3300KV + hélice~3" baseados em extrapolação de curva de eficiência, não em bancada real do par escolhido (TWR$\approx$2,7--3,5; autonomia $\approx$10,5min). & mitigado parcial \\
4 & Tensão exata não garantida em pontos sem conector acima de 100V. & \statushaz \\
5 & Interferência eletromagnética do próprio circuito medido no link WiFi. & \statushaz \\
6 & Contato físico em até 220V exige isolamento rigoroso. & \statusok \\
7 & Bateria danificada em queda perto de circuito energizado — risco composto. & \statusok \\
8 & Falha compartilhada de bateria de eletrônica derrubando os dois ESP32. & \statusok \\
9 & Falha compartilhada de câmera entre ESP32~\#1 e \#2. & \statusok \\
10 & Ausência de plano de terra (PCB 1 camada) aumenta suscetibilidade a ruído. & mitigado parcial \\
11 & Precisão do encaixe do plugue incompatível com tolerância de pilotagem manual. & \statusok \\
12 & Drone preso ao alvo indefinidamente se o link cair durante contato. & \statusok \\
\end{longtable}

% =====================================================
% 6. BOM
% =====================================================
\section{Levantamento de custos (BOM)}

Os valores abaixo combinam cotações diretas de mercado brasileiro (agosto/2026, marcadas com \pesq) e estimativas de engenharia para itens sem padronização de preço (componentes equivalentes, faixa de mercado observada). Colunas \textbf{Mín.}, \textbf{Méd.} e \textbf{Máx.} representam o cenário mais barato, o cenário médio esperado e o cenário mais caro plausível por item.

\subsection{Estrutura e propulsão}

\begin{longtable}{@{}p{5.4cm}p{0.9cm}p{1.7cm}p{1.7cm}p{1.7cm}p{2.1cm}@{}}
\toprule
\textbf{Item} & \textbf{Qtd} & \textbf{Mín.\ un.} & \textbf{Méd.\ un.} & \textbf{Máx.\ un.} & \textbf{Total (méd.)} \\
\midrule
\endhead
\bottomrule
\endfoot
Motor brushless RS1606 3300KV & 6 & R\$~35,00 & R\$~50,00 & R\$~70,00 & R\$~300,00 \\
ESC brushless 30A (dedicado por motor) & 6 & R\$~18,00 & R\$~25,00 & R\$~35,00 & R\$~150,00 \\
Kit hélices 3" (jogo + reservas) & 1 & R\$~20,00 & R\$~35,00 & R\$~50,00 & R\$~35,00 \\
Filamento PLA/PETG (chassi) & 1 & R\$~30,00 & R\$~45,00 & R\$~60,00 & R\$~45,00 \\
\midrule
\multicolumn{5}{r}{\textbf{Subtotal — Mín.\ R\$~368,00 \quad Méd.\ R\$~530,00 \quad Máx.\ R\$~740,00}} & \\
\end{longtable}

\subsection{Controle e comunicação}

\begin{longtable}{@{}p{5.4cm}p{0.9cm}p{1.7cm}p{1.7cm}p{1.7cm}p{2.1cm}@{}}
\toprule
\textbf{Item} & \textbf{Qtd} & \textbf{Mín.\ un.} & \textbf{Méd.\ un.} & \textbf{Máx.\ un.} & \textbf{Total (méd.)} \\
\midrule
\endhead
\bottomrule
\endfoot
ESP32-S3 (módulo dev)\pesq & 2 & R\$~27,34 & R\$~45,00 & R\$~74,80 & R\$~90,00 \\
Módulo LoRa SX1276/RFM95\pesq & 2 & R\$~70,00 & R\$~79,90 & R\$~90,00 & R\$~159,80 \\
\midrule
\multicolumn{5}{r}{\textbf{Subtotal — Mín.\ R\$~194,68 \quad Méd.\ R\$~249,80 \quad Máx.\ R\$~329,60}} & \\
\end{longtable}

\subsection{Câmeras}

\begin{longtable}{@{}p{5.4cm}p{0.9cm}p{1.7cm}p{1.7cm}p{1.7cm}p{2.1cm}@{}}
\toprule
\textbf{Item} & \textbf{Qtd} & \textbf{Mín.\ un.} & \textbf{Méd.\ un.} & \textbf{Máx.\ un.} & \textbf{Total (méd.)} \\
\midrule
\endhead
\bottomrule
\endfoot
Câmera térmica MLX90640\pesq & 1 & R\$~180,00 & R\$~260,00 & R\$~400,00 & R\$~260,00 \\
Câmera FPV + VTX (navegação) & 1 & R\$~90,00 & R\$~150,00 & R\$~250,00 & R\$~150,00 \\
Câmera dedicada ESP32~\#2 (OV2640)\pesq & 1 & R\$~25,00 & R\$~40,00 & R\$~55,00 & R\$~40,00 \\
\midrule
\multicolumn{5}{r}{\textbf{Subtotal — Mín.\ R\$~295,00 \quad Méd.\ R\$~450,00 \quad Máx.\ R\$~705,00}} & \\
\end{longtable}

\subsection{Sensores de medição}

\begin{longtable}{@{}p{5.4cm}p{0.9cm}p{1.7cm}p{1.7cm}p{1.7cm}p{2.1cm}@{}}
\toprule
\textbf{Item} & \textbf{Qtd} & \textbf{Mín.\ un.} & \textbf{Méd.\ un.} & \textbf{Máx.\ un.} & \textbf{Total (méd.)} \\
\midrule
\endhead
\bottomrule
\endfoot
Sensor de campo EMF (não-contato) & 1 & R\$~30,00 & R\$~50,00 & R\$~70,00 & R\$~50,00 \\
Clamp de corrente SCT-013\pesq & 1 & R\$~47,50 & R\$~55,00 & R\$~63,10 & R\$~55,00 \\
Sensor Hall/shunt de contato (ACS712) & 1 & R\$~15,00 & R\$~20,00 & R\$~28,00 & R\$~20,00 \\
Optoacoplador PC817 (módulo)\pesq & 2 & R\$~7,00 & R\$~8,00 & R\$~12,00 & R\$~16,00 \\
Fusível/PTC + proteção & 1 & R\$~5,00 & R\$~10,00 & R\$~15,00 & R\$~10,00 \\
\midrule
\multicolumn{5}{r}{\textbf{Subtotal — Mín.\ R\$~111,50 \quad Méd.\ R\$~151,00 \quad Máx.\ R\$~200,10}} & \\
\end{longtable}

\subsection{Efetuadores}

\begin{longtable}{@{}p{5.4cm}p{0.9cm}p{1.7cm}p{1.7cm}p{1.7cm}p{2.1cm}@{}}
\toprule
\textbf{Item} & \textbf{Qtd} & \textbf{Mín.\ un.} & \textbf{Méd.\ un.} & \textbf{Máx.\ un.} & \textbf{Total (méd.)} \\
\midrule
\endhead
\bottomrule
\endfoot
Micro servo (SG90 ou equiv.)\pesq & 3 & R\$~16,90 & R\$~19,00 & R\$~25,00 & R\$~57,00 \\
Plugue NBR~14136 + adaptadores & 1 & R\$~15,00 & R\$~25,00 & R\$~35,00 & R\$~25,00 \\
Funil de captura + filamento extra & 1 & R\$~10,00 & R\$~15,00 & R\$~20,00 & R\$~15,00 \\
Hardware das garras (dobradiças, fixação) & 1 & R\$~20,00 & R\$~35,00 & R\$~50,00 & R\$~35,00 \\
\midrule
\multicolumn{5}{r}{\textbf{Subtotal — Mín.\ R\$~95,70 \quad Méd.\ R\$~132,00 \quad Máx.\ R\$~180,00}} & \\
\end{longtable}

\subsection{Energia}

\begin{longtable}{@{}p{5.4cm}p{0.9cm}p{1.7cm}p{1.7cm}p{1.7cm}p{2.1cm}@{}}
\toprule
\textbf{Item} & \textbf{Qtd} & \textbf{Mín.\ un.} & \textbf{Méd.\ un.} & \textbf{Máx.\ un.} & \textbf{Total (méd.)} \\
\midrule
\endhead
\bottomrule
\endfoot
Célula Li-ion 18650/21700 alta descarga\pesq & 4 & R\$~25,00 & R\$~35,00 & R\$~45,00 & R\$~140,00 \\
BMS (pack propulsão) & 1 & R\$~25,00 & R\$~40,00 & R\$~60,00 & R\$~40,00 \\
Bateria eletrônica (domínio medição) & 1 & R\$~20,00 & R\$~30,00 & R\$~40,00 & R\$~30,00 \\
Bateria backup ESP32~\#2 & 1 & R\$~10,00 & R\$~17,00 & R\$~25,00 & R\$~17,00 \\
\midrule
\multicolumn{5}{r}{\textbf{Subtotal — Mín.\ R\$~155,00 \quad Méd.\ R\$~227,00 \quad Máx.\ R\$~305,00}} & \\
\end{longtable}

\subsection{Iluminação, PCB e montagem}

\begin{longtable}{@{}p{5.4cm}p{0.9cm}p{1.7cm}p{1.7cm}p{1.7cm}p{2.1cm}@{}}
\toprule
\textbf{Item} & \textbf{Qtd} & \textbf{Mín.\ un.} & \textbf{Méd.\ un.} & \textbf{Máx.\ un.} & \textbf{Total (méd.)} \\
\midrule
\endhead
\bottomrule
\endfoot
LED difuso + driver/resistor & 1 & R\$~10,00 & R\$~17,00 & R\$~25,00 & R\$~17,00 \\
Placa perfurada (protoboard) & 6 & R\$~8,00 & R\$~12,00 & R\$~15,00 & R\$~72,00 \\
Fios, conectores, solda, hardware geral & 1 & R\$~50,00 & R\$~75,00 & R\$~100,00 & R\$~75,00 \\
\midrule
\multicolumn{5}{r}{\textbf{Subtotal — Mín.\ R\$~108,00 \quad Méd.\ R\$~164,00 \quad Máx.\ R\$~215,00}} & \\
\end{longtable}

\subsection{Resumo geral e comparação com o orçamento}

\begin{longtable}{@{}p{6.5cm}p{2.3cm}p{2.3cm}p{2.3cm}@{}}
\toprule
\textbf{Categoria} & \textbf{Mín.} & \textbf{Méd.} & \textbf{Máx.} \\
\midrule
\endhead
\bottomrule
\endfoot
Estrutura e propulsão & R\$~368,00 & R\$~530,00 & R\$~740,00 \\
Controle e comunicação & R\$~194,68 & R\$~249,80 & R\$~329,60 \\
Câmeras & R\$~295,00 & R\$~450,00 & R\$~705,00 \\
Sensores de medição & R\$~111,50 & R\$~151,00 & R\$~200,10 \\
Efetuadores & R\$~95,70 & R\$~132,00 & R\$~180,00 \\
Energia & R\$~155,00 & R\$~227,00 & R\$~305,00 \\
Iluminação, PCB e montagem & R\$~108,00 & R\$~164,00 & R\$~215,00 \\
\midrule
\textbf{TOTAL DO PROJETO} & \textbf{R\$~1.327,88} & \textbf{R\$~1.903,80} & \textbf{R\$~2.674,70} \\
\end{longtable}

\destaque{Orçamento cabe no cenário provável, com folga reduzida:}{com o teto revisado para R\$~2.000,00 e a impressão 3D fora do orçamento (fornecida pela faculdade), tanto o cenário mínimo (R\$~1.327,88) quanto o médio (R\$~1.903,80) ficam dentro do teto, com folga de R\$~672,12 e R\$~96,20 respectivamente. A troca para motor brushless (Rev.~1.2) reduziu a folga do cenário médio de R\$~312,20 para R\$~96,20 — margem bem mais apertada que antes.}

\aviso{Risco remanescente — cenário máximo:}{o cenário mais caro (R\$~2.674,70) ainda ultrapassa o teto de R\$~2.000,00 em R\$~674,70 (era R\$~374,70 antes da troca de motor). O conjunto motor+ESC brushless (R\$~312 de oscilação combinada) passou a rivalizar com a câmera térmica MLX90640 (R\$~180 a R\$~400, variação de R\$~220) como maior fator de risco orçamentário, seguidos pela câmera FPV+VTX e pelo pack de energia. Cotação formal do motor/ESC antes da compra é recomendada dado o novo impacto no teto.}

% =====================================================
% 7. DECISÕES
% =====================================================
\section{Registro de decisões}

\begin{longtable}{@{}p{4.4cm}p{4.4cm}p{5.4cm}@{}}
\toprule
\textbf{Decisão} & \textbf{Alternativa considerada} & \textbf{Justificativa} \\
\midrule
\endhead
\bottomrule
\endfoot
Hexacóptero & Quadricóptero / octocóptero & Meio-termo entre custo e redundância de motor \\
Motor brushless RS1606~3300KV + hélice~3" & Motor de escova (decisão original, Rev.~1.1); motor brushless 2205/2207 & Verificação de TWR (peso pessimista $\approx$991g) mostrou que motor de escova não sustenta o peso; motor 2205/2207 sustenta com folga mas puxa corrente ($>$20A/motor) incompatível com a bateria do BOM; RS1606+hélice~3" fecha TWR seguro ($\approx$3--4$\times$) e autonomia $\approx$11,7min (RNF05) \\
Manter a bateria de propulsão original (4$\times$~21700 alta descarga, 4,2Ah) & Redimensionar a bateria para a nova autonomia & Evita retrabalho na célula já escolhida; hélice~3" (menos eficiente que 4") absorve o excesso de autonomia e traz o tempo de voo para dentro de 8--12min sem mexer na bateria \\
Frame em 3 peças modulares (hub + braço~$\times$6 + bandeja de eletrônica), encaixe lingueta/rasgo + parafuso M3 & Frame em peça única & Peça única (330$\times$290mm) não cabe na mesa de impressão de 220$\times$220mm disponível; encaixe lingueta/rasgo resiste a flexão melhor que topo-a-topo só parafusado \\
Bandeja de eletrônica separada, aparafusada no hub por standoffs & ESP32/LoRa direto no hub & Hub não tem área livre para 2$\times$ESP32-S3 (70$\times$28mm) + 2$\times$LoRa sem competir com o favo de mel e os braços; bandeja em segundo andar resolve o conflito de espaço \\
Base de servo (SG90) presente nos 6 braços, não só nos 3 efetuadores & Base só nos 3 braços efetuadores & Mantém a peça do braço única/intercambiável para impressão; custo de peso extra nos 3 braços sem efetuador é desprezível \\
Abordagem manual via FPV & Navegação autônoma completa & Reduz escopo de software; drone é secundário à medição \\
Três efetuadores especializados & Um único efetuador genérico & Cada modo de medição tem exigência física distinta \\
Triagem autônoma às cegas & Operador pré-define o modo & Exigência do projeto: sistema deve classificar sozinho \\
Li-ion 18650/21700 na propulsão & LiPo tradicional & Menor risco de incêndio em queda perto de circuito energizado \\
Dois domínios de bateria & Bateria única com BEC & Evita brownout do ESP32 durante o contato de medição \\
Dois ESP32 com papéis distintos & Um único ESP32 fazendo tudo & Redundância de rádio e supervisor de segurança independente \\
Navegação de emergência por marcador visual & Dead-reckoning puro & Mais robusto contra deriva; problema de visão contido a um alvo fixo \\
Placa de medição em duas metades separadas & Duas zonas na mesma placa & Isolamento físico real, independente da solda manual \\
PCB em protoboard (1 camada) & Fabricação profissional sob encomenda & Economia de orçamento e prazo \\
\end{longtable}

% =====================================================
% 8. PENDÊNCIAS
% =====================================================
\section{Pendências e trabalho futuro}

\begin{itemize}[itemsep=3pt]
  \item Validação em bancada real do motor RS1606~3300KV com hélice de 3" (empuxo e corrente em throttle parcial) — a autonomia estimada ($\approx$10,5min) usa extrapolação da curva de eficiência de outro motor, não dado medido do par escolhido.
  \item Confirmação do padrão de furação real do RS1606 (assumido 12$\times$12mm, 4$\times$M2 — padrão típico da classe 1606, não confirmado contra datasheet do fabricante) antes de imprimir o braço em lote.
  \item Confirmação do footprint real do módulo LoRa (SX1276/RFM95) para apertar o encaixe da bandeja de eletrônica — hoje é uma baia genérica com abraçadeira por falta de dado de fixação confirmado.
  \item Cotação formal do motor brushless RS1606~3300KV e do ESC dedicado — juntos, tornaram-se o maior fator de oscilação do orçamento (R\$~312 de variação combinada), reduzindo a folga do cenário médio para R\$~96,20.
  \item Peso da estrutura impressa (hub + 6 braços + bandeja) já confirmado por CAD real: 264g. Falta confirmar o peso real dos demais componentes (motor, ESC, bateria, eletrônica, sensores) contra a estimativa pessimista usada na verificação de empuxo ($\approx$1.110g no total).
  \item Fechamento do BOM com cotações formais (não estimativas) para os itens sem preço pesquisado.
  \item Cotação formal da câmera térmica MLX90640 antes das demais compras — é o item de maior risco de estourar o teto de R\$~2.000,00 no cenário máximo.
  \item Escolha de modelos específicos de sensor (clamp de corrente, sensor de campo EMF, sensor Hall/shunt de contato).
  \item Projeto das carcaças dos 3 efetuadores (garra de campo, garra de contato, funil/guia de captura do plugue NBR~14136) — a base de fixação no braço já existe, falta a peça específica de cada um.
  \item Definição empírica dos timeouts do failsafe (janela curta de reconexão vs.\ janela longa em estado engajado).
  \item Definição do protocolo exato de mensagens LoRa (formato de pacote, taxa de amostragem da telemetria).
  \item Projeto detalhado do aplicativo desktop (stack tecnológica, banco de dados local, layout do dashboard).
\end{itemize}

% =====================================================
% REGISTRO DE APROVAÇÃO
% =====================================================
\section*{Registro de aprovação}
\addcontentsline{toc}{section}{Registro de aprovação}

\begin{longtable}{@{}p{4.2cm}p{4.2cm}@{}}
\textbf{Elaborado por:} & \rule{4cm}{0.4pt} \\[1.4em]
\textbf{Data:} & \rule{4cm}{0.4pt} \\[1.4em]
\textbf{Orientador:} & \rule{4cm}{0.4pt} \\[1.4em]
\textbf{Assinatura:} & \rule{4cm}{0.4pt} \\
\end{longtable}

\vfill
\begin{center}
{\footnotesize\ttfamily\color{voltgray} SkyVolt \;\textbullet\; Documento de Requisitos, Arquitetura \& BOM \;\textbullet\; Revisão 1.3 \;\textbullet\; 14/08/2026}
\end{center}

\end{document}
